\section{Design specifications and tuning}

\subsection{Time-domain specifications}

\subsubsection{Steady-state error}
\[
e_{\mathrm{ss}}=\lim_{t\to\infty}e(t)=\lim_{s\to 0}sE(s)
\]
Steady-state error is computed with the final value theorem when the closed loop is stable.

\subsubsection{Unit-feedback step error}
\[
e_{\mathrm{ss}}=\lim_{s\to 0}\frac{1}{1+L(s)}
\]
For a unit step reference and unity feedback, this formula links low-frequency loop gain to tracking error.

\subsubsection{System type}
\[
L(s)=\frac{K\prod_i(s-z_i)}{s^q\prod_k(s-p_k)},\qquad q=\text{number of integrators}
\]
The system type is the number of poles at the origin in the loop transfer function.

\subsubsection{Step, ramp, and parabolic error constants}
\[
K_p=\lim_{s\to0}L(s),\qquad
K_v=\lim_{s\to0}sL(s),\qquad
K_a=\lim_{s\to0}s^2L(s)
\]
These constants determine steady-state error for polynomial references.

\subsubsection{Reference error table}
\[
\begin{array}{c|ccc}
\text{Type} & \text{step} & \text{ramp} & \text{parabolic}\\ \hline
0 & 1/(1+K_p) & \infty & \infty\\
1 & 0 & 1/K_v & \infty\\
2 & 0 & 0 & 1/K_a
\end{array}
\]
Integral action increases system type and improves low-frequency tracking.

\subsection{Frequency-domain specifications}

\subsubsection{Bandwidth and speed}
\[
\omega_c\uparrow \quad\Longrightarrow\quad t_r\downarrow
\]
Higher crossover frequency usually gives faster response, limited by phase margin, noise, saturation, and unmodelled dynamics.

\subsubsection{Phase margin and damping}
\[
\gamma_M\uparrow \quad\Longrightarrow\quad M_p\downarrow
\]
Larger phase margin generally means more damping and less overshoot.

\subsubsection{P-controller stability limit}
\[
|K_PG(j\bar{\omega})|=1,\qquad \angle G(j\bar{\omega})=-180^\circ
\]
The critical proportional gain is found at the frequency where the plant phase reaches \(-180^\circ\).

\subsection{Ziegler-Nichols and hand tuning}

\subsubsection{Ultimate gain and period}
\[
K_u=\text{gain at sustained oscillation},\qquad T_u=\text{oscillation period}
\]
The Ziegler-Nichols method uses the stability boundary observed from closed-loop oscillation.

\subsubsection{Ziegler-Nichols closed-loop rules}
\[
\begin{array}{c|ccc}
\text{Controller} & K_P & \tau_i & \tau_d\\ \hline
P & 0.5K_u & - & -\\
PI & 0.45K_u & T_u/1.2 & -\\
PID & 0.6K_u & T_u/2 & T_u/8
\end{array}
\]
These rules give heuristic initial controller parameters for subsequent refinement. % TODO: verify from source

\subsubsection{Magnitude conversion}
\[
M_{\mathrm{dB}}=20\log_{10}M,\qquad M=10^{M_{\mathrm{dB}}/20}
\]
This conversion is used when selecting \(K_P\) from a Bode magnitude plot.
